\section*{Osnovni pojmovi}

\begin{description}

  \item[Hegemonija] Politički i civilizacijski sistem koji povezuje veliki broj ljudskih svetova u Mrežu. Njene institucije obuhvataju Senat, Svestvar i izvršnu vlast koju u vreme hodočašća predvodi Meina Gledston. Hegemonija održava trgovinu, diplomatiju i vojnu zaštitu između svetova, ali njeno širenje često znači političku, ekonomsku i kulturnu dominaciju nad kolonijama. Hiperion je tek pred početak rata formalno uključen u Hegemoniju.

  \item[Hiperion] Zabačen kolonijalni svet na rubu prostora Hegemonije, poznat po neobičnim ekosistemima, starim ruševinama i Vremenskim grobnicama. Na njemu se nalaze Kits, Grad Pesnika, More trave, Masiv Uzde i Utvrđenje Hronos. Hiperion je dugo ostao izvan glavnih tokova Mreže, ali pred poslednje hodočašće postaje središte sukoba između Hegemonije, Proteranih i sila povezanih sa Vremenskim grobnicama.

  \item[Vremenske grobnice] Skup drevnih građevina u dolini severno od Masiva Uzde. Njihovo poreklo nije poznato. Okružene su antientropijskim poljima i povezane sa neobičnim vremenskim pojavama. Prema onome što hodočasnici znaju u prvom romanu, Grobnice se kreću unazad kroz vreme i približava se trenutak kada će se njihova polja otvoriti. Njihova prava svrha ostaje nepoznata.

  \item[Šrajk] Tajanstveno metalno biće povezano sa Vremenskim grobnicama i kultom koji ga naziva Bogom bola. Opisuje se kao visoko, humanoidno stvorenje sa više ruku, oštrim sečivima i gotovo neuništivim metalnim telom. Pojavljuje se i nestaje na način koji prkosi uobičajenom shvatanju prostora i vremena. Legende ga povezuju sa Drvetom bola, na kojem su njegove žrtve nabijene na metalne trnove. U prvom romanu nije razjašnjeno šta je Šrajk, ko ga je stvorio niti kome služi.

  \item[Proterani] Ljudske zajednice koje su tokom Hedžire napustile glavne tokove Hegemonije i razvile sopstvenu civilizaciju izvan Mreže. Hegemonija ih uglavnom predstavlja kao opasne neprijatelje i vojnu pretnju. Tokom hodočašća njihova flota prilazi Hiperionu, a sukob sa SILOM prerasta u otvoreni rat. Konzulovo iskustvo, međutim, pokazuje da su Proterani složeniji nego što propaganda Hegemonije sugeriše.

  \item[TehnoCentar] Zajednica samostalnih veštačkih inteligencija koje su se odavno izdvojile iz neposredne ljudske kontrole, ali i dalje održavaju veliki deo tehnološke infrastrukture Hegemonije. TehnoCentar obezbeđuje računarske sisteme, predviđanja i tehnologije na kojima počiva Mreža. Bron i Džoni otkrivaju da unutar njega postoje različite frakcije sa suprotstavljenim odnosom prema čovečanstvu i da je Hiperion za njihove modele neobična i teško predvidiva promenljiva.

  \item[Dalekobacači] Portali koji omogućavaju trenutni prelazak sa jednog sveta Mreže na drugi. Zahvaljujući njima, udaljene planete mogu funkcionisati gotovo kao delovi jednog istog društva: ljudi mogu živeti, raditi ili putovati kroz više svetova bez dugog međuzvezdanog putovanja. Svetovi bez dalekobacača, poput Hiperiona tokom većeg dela njegove istorije, ostaju mnogo izolovaniji od ostatka Hegemonije.

  \item[Vremenski dug] Razlika između vremena koje doživi putnik tokom međuzvezdanog putovanja i vremena koje protekne za ljude koji su ostali iza njega. Brodovi sa Hokingovim pogonom mogu prelaziti velike udaljenosti, ali zbog relativističkih posledica putniku prođu meseci ili godine dok u spoljašnjem svetu mogu proteći decenije. Vremenski dug zato razdvaja porodice, ljubavnike i generacije. Najizraženiji primer u prvom romanu jeste odnos Merina Aspika i Siri: za njega između susreta prolazi relativno malo vremena, dok Siri između njegovih povrataka stari godinama.

\end{description}